\documentclass[11pt,a4paper]{article}

\usepackage[utf8]{inputenc}
\usepackage[T2A]{fontenc}
\usepackage[english,russian]{babel}
\usepackage{lmodern}
\usepackage[margin=2.4cm]{geometry}
\usepackage{microtype}
\usepackage{parskip}
\usepackage{titlesec}
\usepackage{enumitem}
\usepackage{xcolor}
\usepackage{fancyvrb}
\usepackage{hyperref}

\emergencystretch=3em
\hyphenpenalty=1000
\sloppy

\hypersetup{
  colorlinks=true,
  linkcolor=black,
  urlcolor=blue!50!black,
  pdftitle={Дипломный проект по обнаружению аномалий в АСУ ТП --- обзор и описание кода},
  pdfauthor={Daniyal Mapeke},
}

\titleformat{\section}{\Large\bfseries}{\thesection}{0.8em}{}
\titleformat{\subsection}{\large\bfseries}{\thesubsection}{0.8em}{}
\titleformat{\subsubsection}{\normalsize\bfseries}{\thesubsubsection}{0.8em}{}

\titlespacing*{\section}{0pt}{1.6em}{0.6em}
\titlespacing*{\subsection}{0pt}{1.2em}{0.4em}

\newcommand{\code}[1]{\texttt{#1}}

\fvset{fontsize=\small, frame=single, framesep=4pt, xleftmargin=8pt, xrightmargin=8pt}

\setlist[itemize]{itemsep=2pt, topsep=4pt}
\setlist[enumerate]{itemsep=2pt, topsep=4pt}

\begin{document}

\begin{titlepage}
  \centering
  \vspace*{4cm}
  {\Huge\bfseries Дипломный проект:\\обнаружение аномалий в АСУ ТП\par}
  \vspace{1cm}
  {\Large Объяснение простым языком\\и обзор структуры кода\par}
  \vspace{3cm}
  {\large Даниял Мапеке\par}
  \vspace{0.5cm}
  {\large Май 2026\par}
  \vfill
  {\small Сопроводительный документ к дипломной работе\\
  <<Межстендовая обобщаемость и пофичевая атрибуция\\
  детекторов аномалий для промышленных систем управления>>.\par}
\end{titlepage}

\tableofcontents
\clearpage

%======================================================================
\part*{Часть I --- Как ловить кибератаки на заводах:\\о чём этот диплом на самом деле}
\addcontentsline{toc}{section}{Часть I --- Как ловить кибератаки на заводах}
%======================================================================

\section{Мир, в котором живёт этот проект}

Большинство людей представляют кибербезопасность как кражу банковской карты или блокировку ноутбука вирусом-шифровальщиком. Но это лишь половина области. Вторая половина --- и именно в ней живёт эта дипломная работа --- касается компьютеров, которые не просто хранят информацию, а \textbf{управляют реальными физическими объектами}: станциями водоочистки, электросетями, нефтепроводами, химическими реакторами, заводскими конвейерами.

На любом современном заводе или электростанции стоят сотни датчиков (давление, температура, расход, положение клапанов, обороты двигателей) и сотни исполнительных устройств (насосы, клапаны, выключатели, нагреватели). Между ними находится слой небольших специализированных компьютеров --- в совокупности они называются \textbf{автоматизированной системой управления технологическим процессом} (АСУ ТП, по-английски \textbf{ICS}). Эти компьютеры читают датчики и решают, что должны делать исполнительные устройства. Задача оператора в диспетчерской --- следить за этим и вмешиваться, если что-то идёт не так.

Когда кто-то атакует АСУ ТП, последствия не абстрактны. Несколько реальных примеров:

\begin{itemize}
  \item \textbf{Stuxnet (2010).} Вредонос физически разрушил центрифуги на иранском заводе по обогащению урана, незаметно заставив их вращаться чуть-чуть неправильно, --- при этом операторы видели на экранах совершенно нормальные показания.
  \item \textbf{Атака на украинскую энергосеть (2015).} Атакующие зашли в управляющий софт трёх региональных распределительных компаний и нажали <<отключить>> на выключателях, оставив без света около 225\,000 потребителей посреди зимы.
  \item \textbf{Маручи-Шир, Австралия (2000).} Обиженный подрядчик с украденной рацией открыл канализационные клапаны, сбросив около миллиона литров сточных вод в парки и реки.
\end{itemize}

В каждом из этих случаев атака была \emph{видна в показаниях датчиков} --- если бы кто-то наблюдал за правильным паттерном. Это и есть базовая посылка всей исследовательской области: \textbf{если научить компьютер тому, как выглядит <<норма>> в данных датчиков завода, он сможет кричать, когда что-то выглядит <<не нормой>>.} Эта техника называется \textbf{обнаружением аномалий}, и она --- центральный объект данной дипломной работы.

\section{Почему это вообще исследовательская задача}

Казалось бы, очевидный подход --- <<построить модель машинного обучения, которая выучит норму и будет помечать всё необычное>> --- должен быть давно решён. Откройте любую научную статью за последние пять лет --- цифры выглядят великолепно. Точность обнаружения 95\%, 97\%, 99\%. Если бы эти цифры были честными, задача была бы закрыта, и компании по безопасности просто купили бы лучшую модель и поставили её везде.

На практике они этого не делают, потому что \textbf{цифры в статьях не выдерживают столкновения с реальностью.} Эта дипломная работа --- о том, чтобы количественно показать, \emph{как именно} они не выдерживают, и дать области честный ответ вместо вводящего в заблуждение заголовка.

Есть три конкретных способа, которыми опубликованные цифры вводят в заблуждение.

\subsection{Ловушка 1 --- Шкала оценок подыгрывает модели}

Представьте, что вы оцениваете охранника по числу секунд, которые он смотрел на нужный монитор во время взлома. Если взлом длился десять минут, а охранник глянул на монитор одну секунду, вы могли бы сказать <<он обнаружил атаку>> и поставить максимальный балл. Примерно так работает самая популярная метрика в этой области --- \textbf{point-adjust F1}. Если модель пометила \emph{хотя бы один} момент во время атаки, ей засчитывается полное обнаружение всей атаки, даже если остальные её сигналы --- бессмыслица.

В 2022 году вышла неловкая для области критика: оказалось, что по этой метрике модель, помечающая случайные точки, набирает больше, чем настоящий детектор. То есть весь рейтинг моделей в области может быть отчасти артефактом измерения.

Более честная метрика, \textbf{eTaPR}, оценивает модели по тому, какую долю атаки они поймали и насколько быстро --- как если бы охранника оценивали по той доле взлома, за которой он действительно следил, с учётом того, насколько быстро он среагировал. Дипломная работа перезапускает все модели под обеими метриками и показывает, как меняется ранжирование.

\subsection{Ловушка 2 --- Модели работают только на своём родном стенде}

Каждая опубликованная модель обучается на одном конкретном тестовом стенде (HAI, SWaT, WADI, Morris --- это публичные датасеты, собранные исследователями) и тестируется на нём же. Это как если бы производитель беспилотного автомобиля заявлял <<99\% безопасности>>, потому что машина безупречно проехала \emph{по одной испытательной трассе}, на которой её обучали. Такой автомобиль никто не купит.

Дипломная работа берёт модели, обученные на корейском энергетическом стенде (HAI), и тестирует их на газопроводном стенде (Morris), и наоборот. Практически никто этого до конца не делал --- отчасти потому, что эти два датасета не имеют общих датчиков: в HAI 79 признаков про водяное охлаждение и турбины, в Morris --- 16 признаков про трафик Modbus в газопроводе. Они даже не говорят на одном языке.

Чтобы их состыковать, в работе вводится \textbf{вектор типов признаков}: вместо того чтобы сравнивать <<датчик 47 в HAI>> с <<датчиком 11 в Morris>>, каждый датчик помечается тем, к какому \emph{типу} он относится --- давление, состояние насоса, уставка, положение клапана, управляющий сигнал --- и оба датасета проецируются в это общее 6-мерное пространство. Тогда модель, обученная распознавать <<странное поведение давления>> в HAI, можно спросить про то же самое и в Morris.

\subsection{Ловушка 3 --- <<Мы поймали атаку>> не говорит, что это была за атака}

На реальном заводе с 500 датчиками сигнал тревоги вида <<где-то что-то не так>> практически бесполезен. Оператору нужно знать, \emph{какой именно} датчик ведёт себя странно, чтобы действовать --- закрыть тот клапан, перезапустить тот насос, изолировать ту подсистему. В работе это называется \textbf{локализацией} (в отличие от простого \textbf{обнаружения}), и проверяется, способна ли каждая модель ещё и указать на виновника.

Как это проверить: авторы HAI любезно разметили, какую физическую подсистему (водяной контур, контроллер турбины и т.\,д.) атаковала каждая атака. Дипломная работа оценивает, насколько <<подозрения>> модели по каждому датчику совпадают с реально атакованной подсистемой. Некоторые модели, отлично справляющиеся с обнаружением, оказываются ужасными в локализации --- этот компромисс раньше никто не измерял количественно.

\section{Что именно было построено}

Чтобы проводить честные эксперименты, нельзя просто скачать шесть разных open-source моделей, написанных в шести разных стилях шестью разными исследователями, и им довериться. Сравнение <<яблок с яблоками>> требует яблок --- поэтому вся кодовая база построена вокруг одного абстрактного интерфейса, который реализует каждая модель:

\begin{itemize}
  \item \code{fit} --- выучить, как выглядит норма.
  \item \code{score} --- по новым данным выдать число <<странности>> для каждого момента времени.
  \item \code{attribute} --- по новым данным выдать число <<вины>> для каждого датчика, чтобы понять, \emph{какой именно} датчик выглядит странно.
\end{itemize}

За этим интерфейсом живут шесть моделей:

\begin{itemize}
  \item Две \textbf{классические базовые} (Isolation Forest и One-Class SVM) --- хорошо изученные методы 2000-х.
  \item Две \textbf{автоэнкодерные} (Dense и LSTM) --- нейронные сети, которые учатся восстанавливать нормальные данные; всё, что они не могут восстановить, подозрительно.
  \item Две \textbf{современные трансформерные} (USAD и TranAD) --- опубликованы в 2020 и 2022 на топовых конференциях, предположительно лучшее, что есть у области.
\end{itemize}

Плюс три метрики оценки, реализованные с нуля с юнит-тестами (чтобы сравнение было доказуемо <<виной>> метрик, а не багом), система защиты от <<утечки скейлера>> (тонкая бухгалтерская ошибка, раздувающая опубликованные результаты --- в работе добавлен ассерт, который падает, если кто-то её совершает), и проекция типов признаков для межстендового переноса.

Есть также \textbf{демонстрационное веб-приложение}, позволяющее не-исследователю загрузить данные с датчиков, выбрать модель и увидеть, что именно она помечает, --- замыкая круг между <<картинками в научной статье>> и <<штукой, на которую можно ткнуть мышкой>>.

\section{Главные находки, оказавшиеся неожиданными}

Три находки выделяются особенно --- и каждая из них из тех, что заставила бы команду безопасности изменить процесс оценки поставщиков.

\paragraph{Находка 1: наивный перенос модели на другой завод вырождается в подбрасывание монетки.}
Если взять модель, обученную на HAI, и применить пороговое значение тревоги, подобранное на HAI, к Morris, модель по сути просто предсказывает <<всё нормально>> или <<всё атака>> в зависимости от того, какой класс чаще встречается в Morris. Опубликованный F1 отражает базовые частоты на новом заводе, а не интеллект модели. \textbf{Перекалибровка порога на выборке нормальных данных нового завода исправляет большую часть этой проблемы} --- и тогда современные трансформерные модели наконец вырываются вперёд классических, достигая примерно 0{,}38--0{,}44 по честной метрике eTaPR. Не 0{,}95. И даже не 0{,}7. Но это честное число, и это первый случай, когда чёткое преимущество современных моделей над классическими измерено в реалистичных условиях.

\paragraph{Находка 2: внутризаводская абляция выявляет точечную слепоту.}
Когда работа <<прячет>> датчики одной подсистемы на этапе обучения (скажем, датчики контроллера турбины), а потом тестирует на атаках, направленных именно на эту подсистему, детектор слепнет \emph{только к этим атакам} --- точность на атаках на турбину падает с 0{,}41 до 0{,}03, при этом на остальных остаётся примерно такой же. Это показывает, что то, что выглядит как <<универсальный детектор аномалий>>, на самом деле --- набор подсистемно-специфичных детекторов, у которых случайно общие веса.

\paragraph{Находка 3: худший детектор --- лучший локализатор.}
Простой Dense-автоэнкодер, проигрывающий навороченным трансформерам в чистом обнаружении, --- \emph{единственная} модель в сравнении, чьи пофичевые оценки <<вины>> бьют случайное угадывание на всех проверенных подсистемах. Трансформеры концентрируют своё внимание на одной подсистеме независимо от того, какая на самом деле атакована, --- то есть если бы вы купили <<лучший детектор>> по заголовочной цифре, вы получили бы худший инструмент для подсказки операторам, куда смотреть. Модное <<объяснение через механизм внимания>>, которым взволнована трансформерная литература, в итоге даёт \emph{численно идентичные} атрибуции с простым подходом <<ошибка восстановления>> --- это самостоятельный нулевой результат, заметный сам по себе.

\section{Почему это важно за пределами дипломной работы}

Дипломная работа сознательно не пытается изобрести новую модель. У области уже больше моделей обнаружения аномалий, чем она знает, что с ними делать. Чего области не хватало --- так это \textbf{честного измерения}: способа для команды безопасности, оценивающей поставщика, --- или для рецензента, оценивающего новую статью, --- задать правильные вопросы и получить сравнимые ответы.

Практический вывод --- короткий и конкретный чек-лист:

\begin{enumerate}
  \item \textbf{Не доверяйте одной заголовочной цифре F1.} Просите оценки модели хотя бы по двум разным метрикам, включая событийную.
  \item \textbf{Не разворачивайте модель с порогом, настроенным на её обучающем заводе.} Всегда перекалибровывайте порог на выборке нормальной работы \emph{вашего} объекта перед запуском в работу.
  \item \textbf{Тестируйте обнаружение и локализацию отдельно.} <<Ловит атаки>> и <<подсказывает, на какой датчик смотреть>> --- разные способности, и лучшая в одном модель часто не лучшая в другом.
\end{enumerate}

Эти три рекомендации звучат очевидно, когда они сформулированы. Вклад дипломной работы в том, чтобы --- с цифрами, на реальных датасетах, с воспроизводимым кодом --- показать, что текущая исследовательская литература во многом игнорирует все три, и что игнорирование их раздувает заявленную производительность на величину, которая в ряде случаев --- разница между полезным детектором и подбрасыванием монетки.

Если коротко, в одной фразе: \textbf{это измерительный проект, который показывает, какая доля заявленного прогресса в области реальна, а какая --- артефакт способа измерения, и публикует код, чтобы любой желающий мог перепроверить это на своих собственных системах.}

\clearpage

%======================================================================
\part*{Часть II --- Обзор кодовой базы}
\addcontentsline{toc}{section}{Часть II --- Обзор кодовой базы}
%======================================================================

Репозиторий разделён на \textbf{семь верхнеуровневых директорий, каждая из которых играет чёткую роль}, и весь дизайн подчинён одному правилу: \emph{каждый эксперимент должен быть воспроизводим из конфиг-файла, без скрытого состояния.} В ноутбуках не происходит ничего важного --- они только рисуют графики.

\section{Корень репозитория}

\begin{Verbatim}
DiplomaDaniyal/
    CLAUDE.md                 постоянные инструкции для ИИ-ассистента
    PROJECT_PLAN.md           высокоуровневая дорожная карта (фазы, дедлайны, решения)
    README.md                 краткое введение для людей
    pyproject.toml            метаданные Python-проекта (ruff, pytest, mypy)
    requirements.txt          зафиксированные зависимости
    .pre-commit-config.yaml   авто-форматтеры и линтеры на каждый коммит
    .gitignore                что НЕ коммитить (датасеты, чекпойнты, кэши)
    src/                      собственно исследовательский код
    experiments/              конфиг-управляемые раннеры экспериментов
    tests/                    тесты pytest
    notebooks/                тонкие Jupyter-ноутбуки только для презентации
    app/                      демонстрационное веб-приложение
    scripts/                  одноразовые утилиты
    data/                     датасеты (тяжёлая часть --- в .gitignore)
    results/                  всё, что производят эксперименты
    thesis/                   LaTeX-исходник самой дипломной работы
    presentation/             слайды защиты и сопроводительные документы
    tools/                    утилиты обслуживания репозитория
\end{Verbatim}

Два самых важных файла в корне --- \code{CLAUDE.md} (конституция проекта: правила дизайна, границы scope, что считается <<готовым>>) и \code{pyproject.toml} (он делает проект устанавливаемым через pip и в одном месте настраивает все инструменты разработчика). \code{.pre-commit-config.yaml} --- то, что автоматически приводит код к стилю; <<закоммитить уродливый код>> случайно не получится.

\section{\texttt{src/} --- исследовательский код}

Здесь живёт вся логика. Кардинальное правило: \textbf{если число попадает в дипломную работу, функция, которая его произвела, лежит в \code{src/}, имеет тест и вызывалась из конфиг-управляемого эксперимента.} Ноутбуки только рисуют графики.

\begin{Verbatim}
src/
    __init__.py
    config.py             датакласс, описывающий один эксперимент
    data_loader.py        загрузка HAI и Morris в общую схему
    preprocessing.py      масштабирование, скользящие окна, разбиения
    utils.py              сиды, пути, логирование --- скучные основы
    models/               каждый детектор аномалий (всего 6)
    evaluation/           три метрики оценки
    attribution/          оценка локализации (какой датчик виноват?)
    explain/              SHAP-объяснения для классических моделей
    transfer/             межстендовый перенос (HAI/Morris) и LOPO
    inference/            загрузка обученной модели и оценка новых данных
\end{Verbatim}

\subsection{\texttt{config.py} --- контракт, который соблюдает каждый эксперимент}

Один Python-файл с пятью маленькими датаклассами: \code{DataConfig}, \code{ModelConfig}, \code{TrainConfig}, \code{EvalConfig}, \code{ArtifactConfig}, обёрнутыми в \code{RunConfig}. Вы передаёте \code{RunConfig} путь к YAML-файлу --- получаете полностью типизированный объект. Он также производит детерминированный 12-символьный хеш всей конфигурации, который записывается вместе с результатами --- через несколько месяцев можно посмотреть метрику в \code{results/metrics/summary.parquet} и понять, какой YAML её произвёл.

\subsection{\texttt{data\_loader.py} --- одна схема для всех}

Два датасета (HAI и Morris) приходят в радикально разных форматах: HAI --- папка CSV-файлов с метками времени и 79 колонками датчиков; Morris --- ARFF-файл с 16 признаками и без меток времени. Загрузчик прячет всё это и возвращает \code{DatasetBundle} с фиксированной структурой:

\begin{itemize}
  \item \code{features} --- DataFrame только из числовых колонок (меток здесь \emph{никогда} нет, что предотвращает случайную утечку);
  \item \code{labels} --- 0/1 флаги атаки;
  \item \code{attack\_ids} --- к какому типу атаки относится каждая строка;
  \item \code{split} --- назначения train/val/test, форсируемые в момент загрузки (train и val \emph{только} нормальные строки; атаки могут быть только в test).
\end{itemize}

Размещение этого форсирования в загрузчике, а не в ноутбуках --- сознательное решение: его невозможно забыть.

\subsection{\texttt{preprocessing.py} --- защита от утечки скейлера}

Горстка маленьких функций: min-max скейлер, построитель скользящих окон и помощники для получения порога по перцентилю. Самое интересное --- то, что он \emph{отказывается делать}: внутри есть ассерт, который падает, если попытаться обучить скейлер на данных, содержащих строки атак. Этот один ассерт защищает от причины №1 раздутых опубликованных ICS-результатов.

\subsection{\texttt{src/models/} --- шесть детекторов за одним интерфейсом}

\begin{Verbatim}
models/
    base.py                  AnomalyDetector ABC: fit, score, attribute, save, load
    isolation_forest.py
    ocsvm.py
    autoencoder.py           Dense AE
    lstm_autoencoder.py
    usad.py                  трансформерный SOTA (2020)
    tranad.py                трансформерный SOTA (2022)
    __init__.py              реестр, отображающий строки в классы
\end{Verbatim}

\code{base.py} --- весь архитектурный аргумент кодовой базы в 70 строках. Каждая модель --- будь то 20-летний Isolation Forest или трансформер 2022 года --- реализует те же четыре метода: \code{fit}, \code{score}, \code{attribute}, \code{save}/\code{load}. Раннер эксперимента не знает и не интересуется, с каким классом он работает, --- он просто вызывает эти методы. Именно это превращает <<сменить модель, всё остальное оставить тем же>> в однострочное изменение YAML.

Реестр в \code{\_\_init\_\_.py} позволяет конфигу сказать \code{name: "lstm\_ae"} и получить правильный класс. Добавление седьмой модели сводится к: написать класс, зарегистрировать строку в \code{\_\_init\_\_.py}, добавить YAML-конфиг. Больше ничего не меняется.

\subsection{\texttt{src/evaluation/} --- три метрики с юнит-тестами}

\begin{Verbatim}
evaluation/
    pointwise.py        обычные precision/recall/F1 плюс поиск порога по best-F1
    point_adjust.py     PA-F1 (метрика, которой область злоупотребляет)
    etapr.py            событийная метрика (честная)
\end{Verbatim}

Каждая метрика реализована из примитивов, а не импортирована из библиотеки --- ведь сравнение метрик и есть основная задача дипломной работы, и вы не хотите, чтобы баг вендора влиял на результат. У каждой есть ручные юнит-тесты в \code{tests/test\_evaluation.py} с разобранными примерами, чтобы рецензент мог напрямую проверить реализации.

\subsection{\texttt{src/transfer/} --- механизм межстендового переноса}

\begin{Verbatim}
transfer/
    schema_align.py     проекция HAI(79) и Morris(16) через 6-мерный вектор типов
    lopo.py             вспомогательные функции leave-one-process-out внутри HAI
\end{Verbatim}

В \code{schema\_align.py} живёт самая методологически интересная часть: он читает \code{data/feature\_types.yaml} (размеченный вручную <<этот датчик --- давление, тот --- состояние клапана>>) и проецирует любой из двух датасетов в общий 6-мерный вектор типов. \code{lopo.py} поддерживает <<скрыть процесс P1/P2/P3/P4 во время обучения>> --- именно эту операцию использует внутри-HAI абляция.

\subsection{\texttt{src/attribution/} и \texttt{src/explain/}}

\code{attribution/evaluation.py} вычисляет precision@k <<обвинений>> модели по датчикам относительно разметки HAI о том, какой процесс был атакован, --- метрика локализации.

\code{explain/} содержит SHAP-бэкенды атрибуции для классических моделей (Isolation Forest и OC-SVM сами по себе не выдают пофичевых оценок; SHAP их оборачивает, чтобы их можно было сравнивать с автоэнкодерами на равных). \code{\_shap\_backends.py} --- тонкий слой совместимости для разных типов моделей SHAP; \code{attribution.py} --- публичная точка входа.

\subsection{\texttt{src/inference/} --- превращение обученной модели в пригодное к использованию}

\begin{Verbatim}
inference/
    artifact.py             сохранение/загрузка модели + скейлера + порога + метаданных
    pipeline.py             применение артефакта к новому DataFrame; чистая функция
    adapters/
        hai.py              разбор свежего HAI CSV
        morris_gas.py       разбор свежего Morris ARFF/CSV
\end{Verbatim}

Концепция \code{ModelArtifact} важна: когда вы сохраняете обученную модель, вы сохраняете \emph{всё, что нужно, чтобы её снова использовать} --- веса модели, скейлер, обученный на тренировочных данных, порог, откалиброванный на валидации, и метаданные о том, какие признаки она ожидает. Загрузка артефакта и оценка нового файла --- один вызов функции: \code{score\_dataframe(artifact, df)}. CLI в \code{scripts/score\_external.py} и веб-приложение в \code{app/routes.py} --- оба тонкие обёртки вокруг этой одной функции.

\section{\texttt{experiments/} --- конфиг-управляемые раннеры}

\begin{Verbatim}
experiments/
    run.py                  одиночный эксперимент на одном датасете
    run_transfer.py         межстендовый (HAI/Morris) эксперимент
    run_lopo.py             leave-one-process-out эксперимент
    run_attribution.py      оценка пофичевой атрибуции
    configs/                около 48 YAML-файлов, по одному на ячейку эксперимента
\end{Verbatim}

Каждой ячейке эксперимента в дипломной работе (каждой строке в каждой таблице результатов) соответствует YAML в \code{configs/}. Именование систематично: \code{baseline\_hai\_lstm\_ae.yaml}, \code{transfer\_morris\_to\_hai\_tranad\_tgtcal.yaml}, \code{lopo\_hai\_P2\_usad.yaml}. Суффикс \code{\_tgtcal} означает <<порог перекалиброван на нормальных данных целевого домена>> --- это и есть вклад C1 в форме YAML.

Чтобы воспроизвести любое число из дипломной работы: выберите YAML и запустите \code{python -m experiments.run <yaml>}. Результат добавится в \code{results/metrics/summary.parquet}, помеченный хешем конфига, сидом и временной меткой. Никакого скрытого состояния, никакой логики только-в-Jupyter, никакого <<надо было там быть>>.

\section{\texttt{tests/} --- то, что делает остальное достоверным}

\begin{Verbatim}
tests/
    test_data_loader.py         загрузчик возвращает правильную форму, нет утечки меток
    test_preprocessing.py       защита от утечки скейлера действительно падает
    test_evaluation.py          все три метрики на ручных примерах
    test_models_smoke.py        каждая модель может fit/score/attribute на малых данных
    test_attribution.py         precision@k и базовые линии lift
    test_explain.py             SHAP-бэкенды выдают вывод правильной формы
    test_transfer.py            проекция схем сохраняет инварианты
    test_persistence.py         тождественность save -> load
    test_inference_pipeline.py  путь artifact -> score
    test_hai_adapter.py         разбор CSV
    test_morris_adapter.py      разбор ARFF
    test_app_smoke.py           веб-приложение запускается и отдаёт главную страницу
    test_demo_datasets.py       встроенные демо-данные корректны
\end{Verbatim}

Паттерн такой: у метрик и загрузчиков --- \emph{настоящие} тесты с разобранными примерами; у моделей --- smoke-тесты (обучить на 100 фейковых строках, проверить, что не падает). Это распределение соответствует тому, где баги действительно навредили бы дипломной работе: сломанная метрика обесценивает все числа в документе; сломанная модель просто громко падает.

\section{\texttt{app/} --- демонстрационное веб-приложение}

\begin{Verbatim}
app/
    main.py                 точка входа FastAPI --- поднимает сервер
    routes.py               HTTP-эндпойнты (/score, /artifacts, /demo-datasets, ...)
    schemas.py              pydantic-модели для форм запросов/ответов
    static/                 фронтенд (vanilla HTML + JS + CSS)
        index.html
        app.js
        style.css
\end{Verbatim}

Это то, что замыкает круг между <<исследовательскими графиками>> и <<штукой, на которую можно ткнуть>>. Пользователь открывает главную страницу, выбирает обученный артефакт (например, LSTM-AE, обученный на HAI), загружает CSV или выбирает встроенный демо-датасет --- и видит пофайловые оценки аномалий с метриками. Фронтенд сознательно простой --- никакого React и шага сборки --- потому что это приложение --- приложение к дипломной работе, а не продукт.

\section{\texttt{scripts/} --- одноразовые утилиты}

\begin{Verbatim}
scripts/
    score_external.py                CLI-версия того, что делает веб-приложение
    build_demo_datasets.py           создаёт встроенные CSV в data/demo/
    build_shap_attribution.py        генерирует parquet с SHAP-результатами
    build_shap_attribution_figure.py создаёт соответствующий график
\end{Verbatim}

Отличие от \code{experiments/}: файл в \code{scripts/} производит \emph{артефакт} (график, встроенный датасет), а не строку в таблице результатов. Файл в \code{experiments/} --- измерение; файл в \code{scripts/} --- шаг сборки.

\section{\texttt{notebooks/} --- тонкие, только для графиков}

\begin{Verbatim}
notebooks/
    01_data_exploration.ipynb
    02_preprocessing_sanity.ipynb
    03_baseline_results.ipynb
    04_sota_results.ipynb
    05_cross_dataset.ipynb           заглавная глава
    06_metric_sensitivity.ipynb
    07_attribution.ipynb
\end{Verbatim}

Каждый ноутбук читает \code{results/metrics/summary.parquet} (или parquet атрибуции), фильтрует нужные строки и производит графики в \code{results/figures/<chapter>/}. Никакого обучения моделей ни в одном ноутбуке --- это сделало бы результаты невоспроизводимыми. Если график в ноутбуке выглядит неправильно, не дебажат ноутбук --- перезапускают соответствующий эксперимент, и график перегенерируется.

\section{\texttt{data/} --- только маленькие, закоммиченные части}

\begin{Verbatim}
data/
    feature_types.yaml      ручная разметка типов датчиков --- критична для переноса
    demo/                   крошечные встроенные образцы для веб-приложения
        hai_test_sample.csv
        morris_gas_test_sample.csv
        manifest.json
    raw/                    в .gitignore --- сюда кладутся оригиналы HAI/Morris
    processed/              в .gitignore --- производные parquet
    README.md               как скачать сырые датасеты
\end{Verbatim}

\code{feature\_types.yaml} мал по размеру, но методологически критичен: это человеческий ввод, делающий возможной проекцию вектора типов. Всё остальное под \code{data/} --- либо крошечные демо-образцы, либо тяжёлые данные в .gitignore.

\section{\texttt{results/} --- всё, что выдают эксперименты}

\begin{Verbatim}
results/
    metrics/
        summary.parquet         главная таблица результатов
                                --- по одной строке на (run, seed, metric)
        attribution.parquet     пофичевые оценки атрибуции
        attribution_shap.parquet
    figures/                    по подпапке на каждую главу дипломной работы
    checkpoints/                в .gitignore --- обученные артефакты моделей
    external/                   выводы из загрузок пользователей в веб-приложении
\end{Verbatim}

\code{summary.parquet} --- единый источник правды. Каждый ноутбук читает из него; каждое число в дипломной работе восходит к строке в нём. Parquet, а не CSV, потому что строки имеют вложенные колонки (метрики по типам атак), и parquet это поддерживает нативно.

\section{\texttt{thesis/} --- сам документ}

\begin{Verbatim}
thesis/
    main.tex                    скелет документа
    preamble.tex                пакеты, форматирование, кастомные команды
    references.bib              цитирования
    build.sh                    однокомандная сборка
    chapters/                   01_introduction.tex ... 08_conclusion.tex
    frontmatter/                титульник, аннотации EN/RU/KK, задание и т.д.
    appendices/                 артефакты воспроизводимости
    build/                      собранные PDF (в .gitignore)
\end{Verbatim}

Дипломная работа цитирует графики прямо из \code{results/figures/}. Скрипт сборки пересобирает PDF; ничто в документе не вставлено вручную из внешних источников.

\section{Общая картина: поток данных}

Если проследить одно число через всю систему, путь такой:

\begin{enumerate}
  \item \textbf{Сырые данные} в \code{data/raw/} читаются \code{src/data\_loader.py} в \code{DatasetBundle}.
  \item \code{src/preprocessing.py} их масштабирует и (опционально) нарезает на скользящие окна.
  \item \textbf{Модель} из \code{src/models/} обучается, затем её просят оценить тестовую выборку.
  \item \code{src/evaluation/} превращает оценки в числа метрик.
  \item \code{experiments/run.py} записывает строку в \code{results/metrics/summary.parquet}, помеченную хешем конфига.
  \item \textbf{Ноутбук} в \code{notebooks/} читает этот parquet и производит график под \code{results/figures/}.
  \item \code{thesis/main.tex} включает график через \code{\textbackslash includegraphics} и цитирует его.
\end{enumerate}

У каждого шага есть тест. У каждого эксперимента --- YAML. У каждого результата --- хеш конфига. Эта дисциплина и есть настоящий вклад кодовой базы --- именно она позволяет рецензенту поверить числам.

\end{document}
